%%%%%%%%%%%%%%%%%%%%%%%%%%%%%%%%%%%%%%%%%%%%%%%%%%%%%%%%%%%%%%%%%%%%%%
% Lớp 10 | Phương pháp tọa độ trong mặt phẳng (𝑂𝑥𝑦) | Ba đường conic trong mặt phẳng tọa độ
% Dạng 0H9?5-3: Bài toán điểm trên elip
% Số câu: 7  (sinh tự động bởi tach_cau.py; từ nay thầy sửa trực tiếp file này)
%%%%%%%%%%%%%%%%%%%%%%%%%%%%%%%%%%%%%%%%%%%%%%%%%%%%%%%%%%%%%%%%%%%%%%

%%%%%%%%%%%%%%%% TRẮC NGHIỆM 4 PHƯƠNG ÁN %%%%%%%%%%%%%%%%

% nguồn: MaHoaNH_TracNghiem_Lop10
\begin{ex}%[0H9H5-3]%[ma:6463d41b]
	 Cho elip $(E)$ có phương trình $\dfrac{x^2}{16}+\dfrac{y^2}{4}=1$. $M$ là một điểm thuộc $(E)$ sao cho $MF_1=MF_2$. Khi đó, tọa độ của điểm $M$ là
	\choice 
		{ $M_1(0;1)$ và $M_2(0;-1)$} 
		{ \True  $M_1(0;2)$ và $M_2(0;-2)$}
		{ $M_1(-4;0)$ và $M_2(4;0)$} 
		{ $M_1(0;4)$  và $M_2(0;-4)$}
	\loigiai{Phương tình chính tắc  của elip $(E)$ là $\dfrac{x^2}{16}+\dfrac{y^2}{4}=1$ nên $a=4$, $b=2$.\\
		Vì $MF_1=MF_2$ nên $M$ thuộc đường trung trực của $F_1F_2$. Mà trung trực của $F_1F_2$ chính là trục $Oy$, do đó $M$ là giao điểm của elip  $(E)$ với trục $Oy$.\\
	Vậy $M_1(0;2)$ và $M_2(0;-2)$.} 
\end{ex}

% nguồn: MaHoaNH_TracNghiem_Lop10
\begin{ex}%[0H9H5-3]%[ma:12548ac2]
	Cho elip $(E)$ có  phương trình $\dfrac{x^2}{25}+\dfrac{y^2}{9}=1$. Đường thẳng $(d)\colon x=-4$ cắt $(E)$ tại hai điểm $M$, $N$. Khi đó
	 \choice
		{  $MN=\dfrac{9}{25}$}
		 {$MN=\dfrac{18}{25}$}
		 {\True $MN=\dfrac{18}{5}$}
		 {$MN=\dfrac{9}{5}$}
	\loigiai{$M(x;y)$ nằm trên $(d)$ nên  $x=-4$, suy ra
		$$\dfrac{(-4)^2}{25}+\dfrac{y^2}{9}=1\Leftrightarrow \dfrac{y^2}{9}=\dfrac{9}{25}\Leftrightarrow y^2=\dfrac{81}{25}\Leftrightarrow\hoac{&  y=-\dfrac{9}{5}\\& y=\dfrac{9}{5}.}$$
		Vậy $M\left(-4;- \dfrac{9}{5}\right)$ và $N\left(-4; \dfrac{9}{5}\right)$. Do đó $MN=\sqrt{(-4+4)^2 +\left(  \dfrac{9}{5}+\dfrac{9}{5}\right)^2}=\dfrac{18}{5}$.
	} 
\end{ex}


%%%%%%%%%%%%%%%% ĐÚNG/SAI %%%%%%%%%%%%%%%%

% nguồn: MaHoaNH_DungSai_Lop10
\begin{ex}%[0H9H5-3]%[ma:ae4096c2]
	Cho  elip $(E)$ có phương trình chính tắc là $\dfrac{x^2}{a^2} + \dfrac{y^2}{b^2} = 1$, biết rằng $(E)$ đi qua điểm $M(2;1)$ và các đỉnh trên trục nhỏ nhìn hai tiêu điểm dưới một góc vuông. Trong các mệnh đề sau, mệnh đề nào đúng, mệnh đề nào sai?
	\choiceTF 
		{\True  $b=c$}
		 {\True $\dfrac{4}{a^2}+\dfrac{1}{b^2}=1$}
		{\True Các tiêu điểm là  $F_1(-\sqrt{3};0)$; $F_2(\sqrt{3};0)$}
		{Tổng các khoảng  cách từ mỗi điểm thuộc elip đến các tiêu điểm bằng $4\sqrt{6}$}
	\loigiai{\immini{ Gọi $B$ là đỉnh trên trục nhỏ, $F_1$, $F_2$ là hai tiêu điểm. \\ 
		Khi đó tam giác  $F_1BF_2$ vuông cân nên $b=c$. Do đó $a^2=b^2+c^2=2b^2$.\\
		Mặt khác $M\in (E)$ nên $\dfrac{4}{a^2} + \dfrac{1}{b^2} = 1$. Từ đó suy ra  $b^2=3$, $a^2=6$.\\
		Vậy $A_1(-\sqrt{6},0)$, $A_2(\sqrt{6},0)$, $B_1(0; -\sqrt{3})$,  $B_2(0;\sqrt{3})$.}
		{\begin{tikzpicture}[line join=round,line cap=round,>=stealth,thick,scale=0.4]
		\tikzset{every  node/.style={scale=0.9}}
		 \draw[->] (-7.1,0)--(7.1,0) node[below left] {$x$};
		\draw[->] (0,-4.1)--(0,4.1)  node[below left] {$y$};
		\draw  (0,0) node [below left] {$O$};
		\coordinate (O) at (0,0); 
		\draw[name path=H] (O) ellipse (5 cm and 3  cm);
		\coordinate  (F1) at (-4,0);
		\coordinate (A1)  at (-5,0);
		 \coordinate (A2) at (5,0);
		\coordinate (F2) at (4,0); 
		\coordinate  (B2) at (0,3);
		 \coordinate (B1) at (0,-3);
		\foreach \i in {O,A1,B1,F1,F2,B2,A2}{\fill[](\i) circle (1.5pt);} 
		\foreach \i in  {F1,F2}{\draw (\i) node[scale=0.9,below]{$\i$};}
		\draw (-5,0) node[below left]{$-a$};  \draw (0,-3) node[below left]{$-b$};
		\draw (5,0) node[above right]{$a$}; \draw (0,3)  node[above right]{$B$};
		\draw  (F1)--(B2)--(F2);
	\end{tikzpicture}
		 }
	} 
\end{ex}


%%%%%%%%%%%%%%%% TRẢ LỜI NGẮN %%%%%%%%%%%%%%%%

% nguồn: MaHoaNH_TraLoiNgan_Lop10
\begin{bt}%[0H9H5-3]%[ma:6ebafaa0]
	Trong mặt phẳng $(Oxy)$, cho elip $(E)\colon \dfrac{x^2}{25}+\dfrac{y^2}{16}=1$ có các tiêu  điểm $F_1$, $F_2$. Lấy hai điểm $A$, $B\in (E)$ sao cho $AF_1+BF_2=6$. Tính tổng $AF_2+BF_1$.
	 \shortans{14}
	\loigiai{ 
		Elip có  dạng $\dfrac{x^2}{a^2}+\dfrac{y^2}{b^2}=1$ nên $\heva{&a^2=25\\&b^2=16}\Leftrightarrow\heva{&a=5\\&b=4.}$\\
		Do $A$, $B\in  (E)$ nên $\heva{&AF_1+AF_2=2a\\&BF_1+BF_2=2a}\Leftrightarrow\heva{&AF_1+AF_2=10\\&BF_1+BF_2=10.}$\\
		Suy  ra $AF_1+AF_2+BF_1+BF_2=20$.\\
		 Mà $AF_1+BF_2=6$ nên $AF_2+BF_1=14$.
	} 
\end{bt}

% nguồn: MaHoaNH_TraLoiNgan_Lop10
\begin{bt}%[0H9H5-3]%[ma:50441217]
	Trong  mặt phẳng tọa độ $Oxy$, phương trình chính tắc của elip đi qua điểm $M(0;3)$ $N\left(3;-\dfrac{12}{5}\right)$ có dạng $\dfrac{x^2}{a^2}+\dfrac{y^2}{b^2}=1$. Các đỉnh của elip này là $A_1$, $A_2$, $B_1$ và $B_2$. Tính $T=A_1A_2+B_1B_2$.
	 \shortans{16}
	 \loigiai{
		Phương trình chính  tắc elip $\dfrac{x^2}{a^2}+\dfrac{y^2}{b^2}=1$.\\
		Ta  có $\heva{&\dfrac{9}{b^2}=1\\&\dfrac{9}{a^2}+\dfrac{144}{25b^2}=1}\Leftrightarrow\heva{&\dfrac{1}{a^2}=\dfrac{1}{25}\\&\dfrac{1}{b^2}=\dfrac{1}{9}}\Leftrightarrow\heva{&a^2=25\\&b^2=9}\Leftrightarrow\heva{&a=5\\&b=3.}$\\
		\begin{itemize}
			\item Phương trình elip $\dfrac{x^2}{25}+\dfrac{y^2}{9}=1$. 
			\item Tọa độ các đỉnh  $A_1(-5;0)$, $A_2(5;0)$, $B_1(0;-3)$ và $B_2(0;3)$.
			\item $T=A_1A_2+B_1B_2=10+6=16$. 
		\end{itemize}
	 }
\end{bt}

% nguồn: MaHoaNH_TraLoiNgan_Lop10
\begin{bt}%[0H9H5-3]%[ma:26bc8d45]
	Trong mặt phẳng tọa độ $Oxy$, cho điểm $M$ chuyển động trên elip  $(E)\colon \dfrac{x^2}{25}+\dfrac{y^2}{16}=1$ và điểm $D(6;0)$. Gọi $a$ và $b$ lần lượt là giá trị nhỏ nhất và giá trị lớn nhất của $DM$. Khi đó $T=2a+b$ bằng?
	\shortans{13} 
	\loigiai{ 
		\begin{center}
			\begin{tikzpicture}[scale=1,font=\footnotesize,line join=round,line cap=round,>=stealth]
					\def\xmin{-3.5}\def\xmax{3.5}\def\ymin{-2.5}\def\ymax{2.5}\pgfmathsetmacro\c{1.5} 
					\draw[->] (\xmin-0.2,0)--(\xmax+0.2,0) node[below]  {\footnotesize $x$};
					 \draw[->] (0,\ymin-0.2)--(0,\ymax+0.2) node[right] {$y$};
					 \draw (0,0) node [below left] {$O$};
					\clip (\xmin,\ymin)  rectangle (\xmax,\ymax);
					 \path
						(-2.5,0) coordinate (A_1)  node[below left]{$-5$}
						(2.5,0) coordinate (A_2) node[below left]{$5$} 
						(0,-2) coordinate (B_1) node[below  left]{$-4$}
						(0,2)  coordinate (B_2) node[above left]{$4$}
						(3,0) coordinate (D)  node[below]{$6$}
						(-\c,0)  coordinate (F_1) node[above]{$-3$}
						(\c,0) coordinate (F_2)  node[above left]{$3$}
						(1.5,1.6) coordinate (M); 
					\draw  [blue]
						 (0,0) ellipse (2.5cm and 2cm);
					 \draw
						(F_1)--(M)--(F_2)--cycle  (D)--(M);
					\fill (M) node[above]{$M$} circle  (1pt);
					\foreach  \x/\g in {D/40,F_1/-90,F_2/-90}
					\fill (\x) circle (1pt) 
						($(\x)+(\g:3mm)$) node{$\x$}; 
				\end{tikzpicture}
		\end{center}
		Ta có $\heva{&DO-OM\le DM\le  DO+OM\\& OM\le 5\\& DO=6}\Rightarrow 6-5\le DM\le 6+5\Rightarrow 1\le DM\le 11$.\\
		Vậy $DM$ đạt giá trị nhỏ nhất  là $a=1$ và đạt giá trị lớn nhất là $b=11$, do đó $T=2a+b=13$.
	 }
\end{bt}


%%%%%%%%%%%%%%%% TỰ LUẬN %%%%%%%%%%%%%%%%

% nguồn: MaHoaNH_TuLuan_Lop10
\begin{bt}%[0H9H5-3]%[ma:f0644289]
	Cho phương trình chính tắc của elip $(E)\colon \dfrac{x^2}{a^2}+\dfrac{y^2}{b^2}=1$, ($a,b>0$). Biết elip có hai đỉnh trên trục nhỏ cùng với hai tiêu điểm tạo thành một hình vuông có diện  tích bằng $32$. Hãy tính $T=2a+b^2$.
	 \loigiai{
	\immini{ 
		Phương trình  chính tắc của elip $(E)\colon \dfrac{x^2}{a^2}+\dfrac{y^2}{b^2}=1$, ($a>b>0$) và $c^2=a^2-b^2$.\\
		Hai đỉnh trên trục nhỏ $B_1(0;-b)$, $B_2(0;b)$ và hai tiêu cự $F_1(-c;0)$,  $F_2(c;0)$ là bốn đỉnh của một hình vuông nên $b=c$.\\
		Mặt  khác, diện tích hình vuông bằng $32$ nên\\
		 $2c\cdot 2b=32\Rightarrow b^2=8$.\\
		Vậy $b=c=8$ và $a^2=b^2+c^2=16\Rightarrow a=4$.}{\begin{tikzpicture}[scale=1,font=\footnotesize,line join=round,line cap=round,>=stealth]
		\def\xmin{-2.5}\def\xmax{2.5}\def\ymin{-2}\def\ymax{2}\pgfmathsetmacro\c{1.4} 
		\draw[->] (\xmin-0.2,0)--(\xmax+0.2,0) node[below] {\footnotesize $x$}; 
		\draw[->] (0,\ymin-0.2)--(0,\ymax+0.2) node[right]  {$y$};
		\draw  (0,0) node [below left] {$O$};
		\clip (\xmin,\ymin) rectangle  (\xmax,\ymax);
		 \path
			 (-2,0) coordinate (A_1)
			(2,0)  coordinate (A_2)
			(0,-1.4) coordinate (B_1) 
			(0,1.4)  coordinate (B_2)
			(-\c,0)  coordinate (F_1)
			 (\c,0) coordinate (F_2);
		 \draw [blue]
			(0,0) ellipse (2cm and  1.4cm);
		 \draw
			(B_1)--(F_1)--(B_2)--(F_2)--cycle; 
		\foreach  \x/\g in {F_1/-90,F_2/-90,B_1/-50,B_2/50}
		\fill  (\x) circle (1.5pt)
			 ($(\x)+(\g:3mm)$) node{$\x$};
	\end{tikzpicture}
		 }
		Phương trình chính tắc của elip là $\dfrac{x^2}{16}+\dfrac{y^2}{8}=1$  và $T=2a+b^2=2\cdot 4+\left(2\sqrt{2}\right)^2=16$.
	 }
\end{bt}
